% Shared section: interpretation and limitations of the wording vs. sampling decomposition
The experiment shows that the question matters for levels: the Cantril ladder yields lower answers than the WVS life-satisfaction question, and this difference fully accounts for the lower average level of Gallup data in high-income countries. But the cross-country pattern of the Gallup--WVS gap, which is what drives differences in correlations with GDP and in country rankings, is not explained by the question. It is explained by the residual, i.e.\ by differences between samples in the broad sense: sampling frames, interview modes \citep{dolan_happy_2016,conti_survey_2011}, questionnaire context \citep{deaton_financial_2012,deaton_understanding_2016}, translations, and time.

Several limitations should be kept in mind. First, the experiment covers ten high-income countries, whereas the gap varies most across income levels. The question effect could be different in low- and middle-income countries, where both Gallup and the WVS mostly interview face-to-face; only an experiment in these countries can settle this. Second, the residual lumps together several factors that the design cannot separate. In particular, the question effect was measured with the survey's own translations, toward the end of a questionnaire on global policies, and online; it may differ from the effect that would be obtained with official Gallup and WVS translations, at the beginning of an interview, and face-to-face or by phone. Third, the decomposition assumes that the question effect measured in 2025 applies to earlier years; for several countries, the latest WVS survey dates from the 2000s. Fourth, with ten countries, cross-country statistics are imprecise, as shown by the bootstrap confidence intervals. Finally, online panel respondents report markedly lower well-being than Gallup and WVS respondents asked the same question, a sample and mode effect that deserves attention in its own right as online panels become common in well-being research.
